% =============================================================================
% PACKAGES.TEX — Paquetes para texto académico de física computacional
% =============================================================================
% Orden de carga importa. Agrupados por función; las dependencias de orden
% están indicadas con (→ debe ir después de X).
% =============================================================================

% ─── Encoding y lenguaje ─────────────────────────────────────────────────────
\usepackage[utf8]{inputenc}
\usepackage[T1]{fontenc}
\usepackage[spanish,es-noshorthands]{babel}
\deactivatequoting
\usepackage{csquotes}              % comillas correctas con babel

% ─── Tipografía ──────────────────────────────────────────────────────────────
\usepackage{microtype}
\usepackage{inconsolata}           % fuente mono para listings

% ─── Matemáticas ─────────────────────────────────────────────────────────────
\usepackage{amsmath, amssymb, amsthm, mathtools}
\usepackage{bm}
\usepackage{physics}               % \abs, \norm, \tr, \grad, \div, \curl,
                                   % \laplacian, \dv, \pdv, \bra, \ket, etc.
\usepackage{cancel}                % cancelar términos en derivaciones
% \usepackage{tensor}              % descomentar si se usa notación indicial

% ─── Unidades y números ──────────────────────────────────────────────────────
\usepackage{siunitx}
\sisetup{
  output-decimal-marker = {,},     % separador decimal = coma (español)
  per-mode              = symbol,  % m/s en lugar de m s^{-1}
  separate-uncertainty  = true,    % ±
  group-minimum-digits  = 4,       % 1000 sin separador, 10 000 con
}

% ─── Tablas ──────────────────────────────────────────────────────────────────
\usepackage{booktabs}
\usepackage{array}
\usepackage{tabularx}
\usepackage{longtable}
\usepackage{multirow}
\usepackage{colortbl}              % \rowcolor, \cellcolor, \arrayrulecolor

% ─── Figuras y flotantes ─────────────────────────────────────────────────────
\usepackage{graphicx}              % \includegraphics (png, jpg, pdf)
\usepackage{caption}
\usepackage{subcaption}            % subfiguras
\usepackage{float}                 % posicionamiento [H]
\usepackage{wrapfig}               % figuras embebidas en texto
\usepackage{placeins}              % \FloatBarrier

% ─── Listas ──────────────────────────────────────────────────────────────────
\usepackage{enumitem}
\usepackage{pifont}

% ─── Código fuente ───────────────────────────────────────────────────────────
\usepackage{listings}              % (configuración en listings-catppuccin.tex)

% ─── Algoritmos ──────────────────────────────────────────────────────────────
\usepackage[ruled, vlined, linesnumbered]{algorithm2e}

% ─── Cajas decorativas ───────────────────────────────────────────────────────
\usepackage[many]{tcolorbox}
\tcbuselibrary{theorems, breakable, skins, hooks}

% ─── TikZ y diagramas ───────────────────────────────────────────────────────
\usepackage{tikz}
\usetikzlibrary{
  arrows.meta,
  positioning,
  shapes.geometric,
  fit,
  backgrounds,
  calc,
  decorations.pathreplacing,
  patterns,
  matrix,
}
\usepackage{pgffor}
\usepackage{pgfplots}
\pgfplotsset{compat=1.18}
\usepackage{tikz-cd}

% ─── Espaciado y layout ─────────────────────────────────────────────────────
\usepackage{parskip}
\usepackage{setspace}
\usepackage{fancyhdr}

% ─── Bibliografía ────────────────────────────────────────────────────────────
\usepackage[numbers]{natbib}

% ─── Notas al pie ────────────────────────────────────────────────────────────
\usepackage[bottom, hang]{footmisc}

% ─── Estructura del documento ────────────────────────────────────────────────
\usepackage{appendix}
\usepackage{pdfpages}

% ─── Programación LaTeX ──────────────────────────────────────────────────────
\usepackage{etoolbox}
\usepackage{xparse}

% ─── Hipervínculos (cargar PENÚLTIMO) ────────────────────────────────────────
\usepackage{hyperref}
\usepackage{bookmark}              % (→ después de hyperref)

% ─── Referencias inteligentes (cargar ÚLTIMO) ────────────────────────────────
\usepackage[spanish,noabbrev]{cleveref}
\crefname{equation}{ec.}{ecs.}
\Crefname{equation}{Ecuación}{Ecuaciones}
\crefname{figure}{fig.}{figs.}
\Crefname{figure}{Figura}{Figuras}
\crefname{table}{tabla}{tablas}
\Crefname{table}{Tabla}{Tablas}
\crefname{chapter}{cap.}{caps.}
\Crefname{chapter}{Capítulo}{Capítulos}
\crefname{section}{sec.}{secs.}
\Crefname{section}{Sección}{Secciones}
\crefname{algorithm}{alg.}{algs.}
\Crefname{algorithm}{Algoritmo}{Algoritmos}
\crefname{appendix}{ap.}{aps.}
\Crefname{appendix}{Apéndice}{Apéndices}